\section{Results}
\label{sec:results}

\subsection{Plug-in gains across samplers}
\label{sec:res-main}

\authornote{Lead with \tableref{tab:main}. Claims allowed \emph{only if the numbers
support them}:
(i) +\method{} improves every modifiable sampler in \xx{} of \xx{} cells (E2: \pxx\,dB,
CI \ci);
(ii) the gain is larger at $R{=}8$ than at $R{=}4$, which the spectral-SNR term
predicts because more of the sampled data lie at high frequency relative to the ACS;
(iii) the remaining gap to supervised E2E-VarNet is \xx\,dB. Explain it as a
difference in training regime and do not hide it.
If a sampler does \emph{not} benefit (DDS is a plausible candidate, since its CG
already regularises), report it and give the reason.}

\begin{table}[t]
\floatconts{tab:main}
{\caption{Raw multi-coil fastMRI, $T{=}50$, volume-mean PSNR (dB) / SSIM on the test
volumes. Each baseline is tuned per cell; \method{} uses one configuration throughout.
$\Delta$: paired volume-level difference with 95\% bootstrap CI. Grey: supervised
reference.}}
{\scriptsize\setlength{\tabcolsep}{2.6pt}
\resizebox{\linewidth}{!}{\begin{tabular}{@{}ll cc cc cc cc@{}}
\toprule
& & \multicolumn{2}{c}{Brain $R{=}4$} & \multicolumn{2}{c}{Brain $R{=}8$}
  & \multicolumn{2}{c}{Knee $R{=}4$} & \multicolumn{2}{c}{Knee $R{=}8$}\\
\cmidrule(lr){3-4}\cmidrule(lr){5-6}\cmidrule(lr){7-8}\cmidrule(lr){9-10}
Sampler & DC & PSNR & SSIM & PSNR & SSIM & PSNR & SSIM & PSNR & SSIM\\
\midrule
\multirow{2}{*}{\ikdc} & isotropic & \xx & \xx & \xx & \xx & \xx & \xx & \xx & \xx\\
 & +\method & \xx & \xx & \xx & \xx & \xx & \xx & \xx & \xx\\
\multirow{2}{*}{\pigdm} & isotropic & \xx & \xx & \xx & \xx & \xx & \xx & \xx & \xx\\
 & +\method & \xx & \xx & \xx & \xx & \xx & \xx & \xx & \xx\\
\multirow{2}{*}{DDS} & isotropic & \xx & \xx & \xx & \xx & \xx & \xx & \xx & \xx\\
 & +\method & \xx & \xx & \xx & \xx & \xx & \xx & \xx & \xx\\
\multirow{2}{*}{DiffPIR} & isotropic & \xx & \xx & \xx & \xx & \xx & \xx & \xx & \xx\\
 & +\method & \xx & \xx & \xx & \xx & \xx & \xx & \xx & \xx\\
DPS & --- & \xx & \xx & \xx & \xx & \xx & \xx & \xx & \xx\\
\midrule
\multicolumn{2}{@{}l}{Mean $\Delta$ (4 samplers)} & \pxx & \pxx & \pxx & \pxx & \pxx & \pxx & \pxx & \pxx\\
\multicolumn{2}{@{}l}{\quad 95\% CI} & \ci & & \ci & & \ci & & \ci & \\
\midrule
\multicolumn{2}{@{}l}{\textcolor{gray}{E2E-VarNet (supervised)}} & \textcolor{gray}{\xx} & \textcolor{gray}{\xx} & \textcolor{gray}{\xx} & \textcolor{gray}{\xx} & \textcolor{gray}{\xx} & \textcolor{gray}{\xx} & \textcolor{gray}{\xx} & \textcolor{gray}{\xx}\\
\bottomrule
\end{tabular}}}
\end{table}

\subsection{Step-robustness and mechanism}
\label{sec:res-steps}

\authornote{\tableref{tab:steps} and \figureref{fig:bands}. This is the paper's central
finding (E1). Three things to establish:
(i) isotropic samplers \emph{lose} \xx\,dB from their best $T$ to $T{=}200$ even when
tuned per $T$;
(ii) +\method{} is flat or monotone in $T$, so one configuration serves all budgets
(a practical, tuning-free claim);
(iii) mechanism: per-band error of $\hx$ along the reverse trajectory
(\figureref{fig:bands}). In the isotropic samplers the error in the outer two bands
rises during the last \xx\% of steps; with \method{} it does not. The calibration
ratio (realised residual over predicted $p_i+\sigma_y^2$) stays near 1 with \method{}
and drifts to \xx{} for the isotropic samplers.
Relate (i) to the noise-overfitting / early-stopping evidence of
\citet{ozturkler2023smrd}, and (iii) to \citet{burns2026finite}.}

\begin{table}[t]
\floatconts{tab:steps}
{\caption{Step-count sweep, brain, $R{=}8$ (E1 cell in bold). Baselines tuned
\emph{per $T$}; \method{} uses a single configuration. Volume-mean PSNR (dB); $\Delta$
with 95\% bootstrap CI.}}
{\scriptsize\setlength{\tabcolsep}{3.2pt}
\resizebox{\linewidth}{!}{\begin{tabular}{@{}r ccc ccc@{}}
\toprule
& \multicolumn{3}{c}{\ikdc} & \multicolumn{3}{c}{\pigdm}\\
\cmidrule(lr){2-4}\cmidrule(lr){5-7}
$T$ & isotropic & +\method & $\Delta$ [CI] & isotropic & +\method & $\Delta$ [CI]\\
\midrule
10  & \xx & \xx & \pxx{} \ci & \xx & \xx & \pxx{} \ci\\
20  & \xx & \xx & \pxx{} \ci & \xx & \xx & \pxx{} \ci\\
50  & \xx & \xx & \pxx{} \ci & \xx & \xx & \pxx{} \ci\\
100 & \xx & \xx & \pxx{} \ci & \xx & \xx & \textbf{\pxx{}} \ci\\
200 & \xx & \xx & \pxx{} \ci & \xx & \xx & \pxx{} \ci\\
\midrule
best$\,\to\,$200 drop & $-$\xx & $-$\xx & & $-$\xx & $-$\xx & \\
\bottomrule
\end{tabular}}}
\end{table}

\begin{figure}[t]
\floatconts{fig:bands}
{\caption{\textbf{Mechanism.} Left: k-space NMSE of $\hx$ in five radial bands versus
reverse step (brain, $R{=}8$, $T{=}100$; median over test volumes) for isotropic
\pigdm{} (solid) and +\method{} (dashed). Right: calibration ratio
$\langle|r_i|^2\rangle_b/(p_b+\sigma_y^2)$ per band; 1 is perfectly calibrated.}}
{\figspec{4.0cm}{Two panels. (L) 5 lines per method, colour = band (inner$\to$outer),
x = step index, y = log NMSE. (R) same layout, y = calibration ratio, dashed line at 1.
Source: per-step diagnostics logged by the sampler (\texttt{return\_diagnostics=True}).}}
\end{figure}

\subsection{Low SNR and real low-field data}
\label{sec:res-snr}

\authornote{\tableref{tab:snr}. Hypothesis from \equationref{eq:gain}: the spectral-SNR
term matters more as $\sigma_y^2$ grows, so the gain should increase as SNR falls.
Show it in the controlled sweep (fastMRI brain plus $\Psi$-correlated white noise) and
in real 0.3\,T data (E3). For M4Raw also report against the averaged reference, to
show the gain is not an artefact of a noisy reference. If the trend is not monotone,
say so.}

\begin{table}[t]
\floatconts{tab:snr}
{\caption{Controlled SNR sweep (fastMRI brain, $R{=}4$, $T{=}50$) and real 0.3\,T
M4Raw ($R{=}4$, $T{=}50$, brain prior without fine-tuning). Best tuned isotropic
sampler vs.\ the same sampler +\method. SNR in units of \citet{kellman2005snr}.}}
{\scriptsize\setlength{\tabcolsep}{4pt}
\begin{tabular}{@{}ll ccc@{}}
\toprule
Data & Condition & isotropic & +\method & $\Delta$ [95\% CI]\\
\midrule
\multirow{4}{*}{fastMRI brain} & native (SNR \xx) & \xx & \xx & \pxx{} \ci\\
 & SNR $\times\frac12$ & \xx & \xx & \pxx{} \ci\\
 & SNR $\times\frac14$ & \xx & \xx & \pxx{} \ci\\
 & SNR $\times\frac18$ & \xx & \xx & \pxx{} \ci\\
\midrule
\multirow{4}{*}{M4Raw 0.3\,T} & T1w & \xx & \xx & \pxx{} \ci\\
 & T2w & \xx & \xx & \pxx{} \ci\\
 & FLAIR & \xx & \xx & \pxx{} \ci\\
 & all (\textbf{E3}); averaged ref. & \xx & \xx & \textbf{\pxx{}} \ci\\
\bottomrule
\end{tabular}}
\end{table}

\subsection{Pathology regions and hallucinations}
\label{sec:res-roi}

\authornote{\tableref{tab:roi} and \figureref{fig:qual}. The clinical question: does
frequency-resolved DC preserve lesion detail, or does it smooth it away? Report the
box-versus-normal-patch contrast. Choose qualitative examples \emph{by rule}: median
and 90th-percentile $\Delta$ volumes, frozen before inspection. Never hand-pick.
If \method{} reduces global error but not box error, report that as a limitation.}

\begin{table}[t]
\floatconts{tab:roi}
{\caption{fastMRI+ pathology boxes vs.\ size-matched normal patches (test volumes,
$R{=}8$, $T{=}50$, \pigdm{} host). NMSE ($\times10^{-2}$, lower is better), HFEN, and
fraction of boxes with a hallucination-map peak above threshold \citep{bhadra2021hallucination}.}}
{\scriptsize\setlength{\tabcolsep}{3.5pt}
\begin{tabular}{@{}ll cc cc cc@{}}
\toprule
& & \multicolumn{2}{c}{NMSE box} & \multicolumn{2}{c}{NMSE normal} & \multicolumn{2}{c}{Halluc.\ rate (\%)}\\
\cmidrule(lr){3-4}\cmidrule(lr){5-6}\cmidrule(lr){7-8}
Anatomy & \#boxes & iso & +\method & iso & +\method & iso & +\method\\
\midrule
Brain & \xx & \xx & \xx & \xx & \xx & \xx & \xx\\
Knee  & \xx & \xx & \xx & \xx & \xx & \xx & \xx\\
\bottomrule
\end{tabular}}
\end{table}

\begin{figure}[t]
\floatconts{fig:qual}
{\caption{Rule-selected examples (median and 90th-percentile $\Delta$ volume): ground
truth, isotropic \pigdm, +\method, and absolute-error maps, with the fastMRI+ box
outlined.}}
{\figspec{3.6cm}{2 rows (brain, knee) $\times$ 5 columns (GT, iso, +ours, $|$err iso$|$,
$|$err ours$|$). Shared colour scale per row. Box overlay in yellow. Zoom inset on the
box. Selection rule and IDs logged in \texttt{outputs/midl27/qual\_manifest.json}.}}
\end{figure}

\subsection{Ablations and compute}
\label{sec:res-ablation}

\authornote{\tableref{tab:ablation}. Use these rows to decide what stays in the method.
(b) vs.\ (a) isolates the spectral posterior variance, our principled core.
(c) and (d) test the frequency gate; (f) replaces it with the schedule of
\citet{tian2026frequency}. (e) tests residual refinement. (g) uses an oracle spectrum
from the fully sampled reference and bounds the cost of estimating $\hat S$.
(h)--(i) test robustness to a misestimated $\sigma_y$. (j) tests the smoothing.
\tableref{tab:compute} supports the ``negligible overhead'' claim with measured
numbers only.}

\begin{table}[t]
\floatconts{tab:ablation}
{\caption{Ablations (brain, $R{=}8$, \pigdm{} host; calibration volumes for (h)--(j),
test volumes otherwise). PSNR (dB) and $\Delta$ vs.\ row (a).}}
{\scriptsize\setlength{\tabcolsep}{4pt}
\begin{tabular}{@{}cl cc cc@{}}
\toprule
& & \multicolumn{2}{c}{$T{=}50$} & \multicolumn{2}{c}{$T{=}100$}\\
\cmidrule(lr){3-4}\cmidrule(lr){5-6}
& Variant & PSNR & $\Delta$ & PSNR & $\Delta$\\
\midrule
(a) & isotropic $\bP_t=r_t^2\mathbf{I}$ (tuned) & \xx & --- & \xx & ---\\
(b) & + spectral posterior variance, \eqref{eq:pt} & \xx & \pxx & \xx & \pxx\\
(c) & (b) + frequency gate & \xx & \pxx & \xx & \pxx\\
(d) & frequency gate only (isotropic $\bP_t$) & \xx & \pxx & \xx & \pxx\\
(e) & (b or c) + residual refinement & \xx & \pxx & \xx & \pxx\\
(f) & (b) + schedule of \citet{tian2026frequency} & \xx & \pxx & \xx & \pxx\\
(g) & (b) with oracle spectrum $S$ & \xx & \pxx & \xx & \pxx\\
(h) & (b) with $\sigma_y\times0.5$ & \xx & \pxx & \xx & \pxx\\
(i) & (b) with $\sigma_y\times2$ & \xx & \pxx & \xx & \pxx\\
(j) & (b) raw radial bins (no power-law fit) & \xx & \pxx & \xx & \pxx\\
\bottomrule
\end{tabular}}
\end{table}

\begin{table}[t]
\floatconts{tab:compute}
{\caption{Cost per slice on one A100 (brain, $R{=}8$, $T{=}50$, batch 1). Overhead is
relative to the same sampler with isotropic DC.}}
{\scriptsize\setlength{\tabcolsep}{4pt}
\begin{tabular}{@{}l c cc c c@{}}
\toprule
Sampler & NFE & time iso (s) & time +\method{} (s) & overhead (\%) & peak mem.\ (GB)\\
\midrule
\ikdc   & \xx & \xx & \xx & \xx & \xx\\
\pigdm  & \xx & \xx & \xx & \xx & \xx\\
DDS     & \xx & \xx & \xx & \xx & \xx\\
DiffPIR & \xx & \xx & \xx & \xx & \xx\\
\bottomrule
\end{tabular}}
\end{table}
